\chapter{Introduzione}

\section{L'importanza dell'accessibilità ai servizi primari}

Le attività che tutti noi svolgiamo quotidianamente dipendono strettamente dalla disponibilità e dalla distribuzione spaziale dei servizi essenziali, promuovendo il benessere fisico e mentale delle comunità \cite{wu2025}. Ad esempio, basti pensare alle città che non dispongono di un ospedale costringendo i cittadini a spostarsi verso altre località per usufruire di un serivizio essenziale per la salute. 

\section{Il concetto di accessibilità: definizioni ed evoluzione teorica}

Nel linguaggio comune, il termine accessibilità indica la caratteristica di un servizio, di una risorsa o di un ambiente di essere fruibile con facilità da parte di qualsiasi tipologia di utente. Tuttavia, quando si cerca di formalizzare il concetto in ambiti più complessi, come ad esempio la pianificazione urbana o dei trasporti, il principio di accessibilità non risulta banale \cite{geurs2004}.\\
Dunque, se volessimo dare una definizione univoca rischieremo di essere riduttivi, poiché il significato stesso di accessibilità può variare e mutare a seconda della prospettiva che si adotta e al contesto di analisi. \\
\textbf{DA RIFARE LA FRASE SOTTO}\\
Una domanda che può sembrare facile ma una definizione può rimanere limitatoria e quasi restrittiva quelle che può assumere a seconda di cosa chiediamo davvero quando chiediamo accessibilità.\\
Per esempio, si potrebbe definire l'accessibilità attraverso la distanza fisica: più un servizio risulta vicino alla propria abitazione, più appare accessibile \cite{benassai2026}. Tuttavia, se cambiassimo prospettiva, potremmo considerare il concetto di accessibilità non solo come vicinanza al singolo servizio, ma anche come la possibilità di variare la scelta in termini di quantità e diverse opportunità che un cittadino dispone entro un determinato raggio a partire da un punto.\cite{battiston2024,geurs2004}. \\

\noindent
Dal punto di vista teorico, gli studi sull'accessibilità spaziale pongono le loro fondamenta sulla celebre formulazione del geografo e cartografo Waldo Tobler, nonché la \textbf{Prima Legge della Geografia di Tobler} \cite{tobler1970}, la quale recita la seguente frase: \textit{"ogni cosa è correlata a ogni altra, ma le cose vicine sono più correlate di quelle lontane"}. Questo principio introduce il concetto di decadimento con la distanza (\textit{distance decay}), ripreso ampamente in letteratura \cite{mazzulla2024}.\\
Da questa base, altri studiosi hanno potuto introdurre altre definizioni di accessibilità. Citiamo, ad esempio, lo scenziato e professore di geografia Hansen \cite{hansen1959}, il quale definisce l'accessibilità come \textit{"il potenziale di opportunità per l'interazione spaziale"}; Dalvi \& Martin \cite{dalvi1976} l'hanno descritta come \textit{"la facilità con cui qualsiasi attività di uso del suolo può essere raggiunta da una località utilizzando un particolare sistema di trasporto"}; infine 
Karst Geurs e Bert van Wee \cite{geurs2004} evidenziano come l'accessibilità sia utilizzato in diversi ambiti come quello nelle politiche pubbliche, nella pianificazione territoriale e urbanistica, rivestendo di fatto un ruolo cruciale ma che al contempo sia spesso frainteso come concetto, privo di una definizione precisa e misurato attraverso indicatori non sempre adeguati.
Approfondendo quest'ultimo concetto, gli autori discutono di quanto sia difficile definire un concetto e misurazione di accessibilità, proponendo la seguente definizione:
\begin{quote}
\textit{``L'accessibilità esprime la misura in cui i sistemi di uso del suolo e di trasporto consentono a (gruppi di) individui di raggiungere attività o destinazioni per mezzo di una (combinazione di) modalità di trasporto''} \cite{geurs2004}.
\end{quote}

\subsection{Distinzioni concettuali: accessibilità attiva, passiva e percepita}

Prima di analizzare le componenti strutturali dell'accessibilità, è necessario definire alcune distinzioni terminologiche e concettuali, derivanti dalla letteratura, per inquadrare correttamente lo studio dei servizi primari \cite{mazzulla2024, benassai2026}:
\begin{itemize}
    \item \textbf{Accessibilità spaziale all'interno del concetto multidimensionale di accesso}: il concetto generale di accesso ai servizi è intrinsecamente multidimensionale e comprende aspetti sia spaziali che non, come la disponibilità dell'offerta (\textit{availability}) e la sostenibilità economica dei prezzi per l'utente (\textit{affordability}) \cite{benassai2026}. 
    Nelle analisi territoriali e su larga scala, dove le informazioni dettagliate sui prezzi o sull'assortimento dei singoli servizi sono spesso indisponibili, l'attenzione si concentra specificamente sulla dimensione dell'\textit{accessibilità spaziale}, intesa come la capacità di superare l'impedenza fisica e temporale che separa la domanda dall'offerta \cite{benassai2026, mazzulla2024}.
    *riscrivere*
    \item \textbf{Accessibilità attiva vs. Accessibilità passiva}: l'\textbf{accessibilità attiva} misura il grado di facilità con cui un utente, partendo da una determinata zona (\textit{i}), può raggiungere le attività e le opportunità distribuite sul territorio.
    D'altra parte, l'\textbf{accessibilità passiva}, spesso associata al concetto di \textit{attrattività del luogo} (\textit{place accessibility}), adotta una prospettiva opposta, spostando l'attenzione dall'individuo alla destinazione o al servizio. Essa misura, infatti, il grado di facilità con cui uno specifico punto di interesse può essere raggiunto dagli utenti provenienti dalle diverse zone dell'area di studio. \cite{mazzulla2024}. 
    *riscrivere*
    \item \textbf{Accessibilità oggettiva (aggregata) vs. Accessibilità percepita}: le metriche tradizionali si basano sull'integrazione di dati spaziali e temporali che sono quantificabili in maniera "oggettiva"(distanze, tempi di viaggio, numero di strutture), dando luogo a modelli non-comportamentali che descrivono l'accessibilità potenziale offerta dal territorio \cite{mazzulla2024, battiston2024}.
     Tuttavia, il comportamento reale degli individui è fortemente influenzato dalle percezioni soggettive \cite{mazzulla2024}. Studi più recenti pongono forte accento nell'\textbf{accessibilità percepita} (\textit{perceived accessibility}), cioè la valutazione soggettiva da parte di un individuo in merito alla facilità con cui può raggiungere le attività desiderate e condurre una vita soddisfacente grazie al sistema dei trasporti \cite{mazzulla2024}.
     Essa dipende da esperienze personali, abitudini, senso di sicurezza e fattori socio-economici, illustrando come in molti contesti si registri un divario (\textit{mismatch}) tra l'accessibilità calcolata oggettivamente su base cartografica e la percezione vissuta dai cittadini \cite{mazzulla2024}.
    *riscrivere*
\end{itemize}
Oltre ai concetti precedenti, è necessario soffermarsi sull'interazione tra tre elementi cardine:
\begin{itemize}
    \item \textbf{Le persone}, con esigente, capacità e risorse socio-economiche differenti
    \item \textbf{Le opportunità}, intese come le attività e le destinazioni che possono interessare i cittadini come servizi sanitari, negozi, scuole, ...
    \item \textbf{Il sistema di mobilità}, quindi la possibilità effettiva di raggiungere determinati servizi, dipendenti dal territorio e dalle infrastrutture di trasporto disponibili.
\end{itemize}
Questi concetti permettono di fornire una motivazione e comprendere come, all'interno di uno stesso contesto urbano, quartieri caratterizzati da tempi di spostamento simili possano offrire servizi e opportunità differenti.

\subsection{Le quattro componenti dell'accessibilità}
Secondo Geurs e van Wee, una misura ideale di accessibilità dovrebbe tener conto di alcuni componenti fondamentali:
\begin{enumerate}
    \item \textbf{Componente di uso del suolo (\textit{Land-use component})}: riflette il sistema territoriale e si articola in:
    \begin{enumerate}
        \item la quantità, qualità e distribuzione spaziale delle opportunità offerte in ogni destinazione (offerta);
        \item la distribuzione spaziale della popolazione e il relativo fabbisogno di servizi (domanda);
        \item il confronto tra domanda e offerta, che può generare competizione per i servizi con minore capacità (posti letto ospedalieri o posti nelle scuole).
    \end{enumerate}
    \item \textbf{Componente dei trasporti (\textit{Transport component})}: descrive il funzionamento del sistema di trasporti, espressa come disutilità per un individuo di coprire la distanza tra l'origine e la destinazione usando uno specifico mezzo di trasporto. Comprende:
     \begin{enumerate}
        \item Tempo complessivo (di viaggio, di attesa)
        \item Costi (fissi e variabili)
        \item sforzo fisico e psicologico (comprende livello di comfort, rischio incidenti, ...)
    \end{enumerate}
        Questa diversità deriva dal confronto tra domanda e offerta, dove l'offerta di infrastrutture include posizione e le caratterisctiche del viaggio ( velocità massima, orari trasporti pubblici)\\
        La domanda invece riguarda sia il trasporto passeggeri che di merci. 
    \item \textbf{Componente temporale (\textit{Temporal Component})}: considera i vincoli di tempo, quindi non solo la disponibilità delle opportunità nelle diverse fasce della giornata (come gli orari di apertura e chiusura dei servizi) ma anche il budget di tempo che i cittadini possono dedicare agli effettivi spostastamenti e attività.
    
    \item \textbf{Componente individuale (\textit{Individual component})}: riflette i bisogni, le capacità e le opportunità che variano a seconda di diversi fattori come età, reddito, livello di istruzione dei diversi individui.\\
        Tali caratteristiche influenzano sia la scelta del mezzo di trasporto (che può variare o meno con il possesso di una patente di guida), sia la pertinenza stessa delle destinazioni raggiungibili.
    \end{enumerate}
Dalla combinazione di questi aspetti, gli autori derivano quattro prospettive di misurazione dell'accessibilità che descrivono come essa debba essere misurata e osservata.
\begin{table}[h]
  \centering
  \begin{tabular}{|l|c|r|}
    \hline
    Prospettiva & Che cosa prende in esame & Esempio \\
    \hline
    Basata sulle infrastrutture & Prestazioni del sistema di trasporto & Velocità media o congestione \\
    Basata sui luoghi &	Opportunità raggiungibili a partire da determinate località	& Numero di opportunità entro un tempo di viaggio \\
    Basata sulle persone &	Attività praticabili da una persona considerando i suoi vincoli	& Compatibilità con impegni obbligatori e tempo disponibile\\
    Basata sull’utilità &	Benefici che le persone ricavano dall’accesso alle attività	& Valutazione economica delle alternative disponibili\\
   \hline
  \end{tabular}
  \caption{Esempio di tabella}
  \label{tab:esempio}
\end{table}

\section{Tassonomia delle metriche di Misurazione}
Sulla base della letteratura di riferimento \cite{geurs2004, mazzulla2024, battiston2024, benassai2026}, è possibile classificare le metriche di accessibilità in quattro grandi famiglie, analizzandone i rispettivi punti di forza e i limiti metodologici.

\subsection{Misure Basate sull'infrastruttura}
Queste metriche mirano l'esclusiva performance della rete di trasporto, utilizzando indicatori come la velocità media di percorrenza o i ritardi dovuti alla congestione \cite{geurs2004}. \\
Un vantaggio nell'uso di queste misure risiede nella facilità di calcolo e nell'immediata comunicabilità. Tuttavia, il loro limite è l'omissione della componenti di uso del suolo (\textit{"Land use"}): un quartiere centrale congestionato risulterebbe di bassa accessibilità nonostante disponga di un'offerta densa di servizi primari raggiungibili a pochi metri di distanza \cite{geurs2004}.

\subsection{Misure basate sulla Localizzazione(\textit{Location-based})}

Queste misure analizzano l'accessibilità a livello aggregato, integrando l'uso dei trasporti e del suolo \cite{geurs2004, mazzulla2024}. Rappresentano la famiglia di indicatori più utilizzata per lo studio dei servizi primari e si suddividono in quattro sottocategorie principali:
\begin{enumerate}
    \item \textbf{Misure di separazione spaziale o di distanza minima (\textit{Closeness/Minimum Distance})}:  misurano la distanza o il tempo di viaggio minimo dalla residenza al servizio primario più vicino \cite{battiston2024,wu2025,benassai2026}.
    Per calcolare i tempi di percorrenza si possono utilizzare diversi algoritmi di cammino minimo, (\textit{shortest-path}, come l'algoritmo di Dijkstra) su gradi stradali (es. OpenstreetMap), oppure mediante superfici di frizione (\textit{friction surfaces}), ovvero griglie raster ad alta risoluzione che assegnano una velocità di attraversamento a ciascuna cella in funzione della rete viaria, della copertura del suolo e della pendenza topografica \cite{wu2025, benassai2026}.
    \begin{itemize}
        \item \textit{Vantaggi}: restituiscono un'unità di misura intuitiva (minuti o metri) e non richiedono la calibrazione di parametri comportamentali \cite{benassai2026}.
        \item \textit{Svantaggi}: ignorano la dimensione dell'offerta (avere un solo servizio o dieci servizi a 5 minuti produce il medesimo risultato) e assumono implicitamente che l'utente scelga sempre e solo il servizio più vicino \cite{tariverdi2023, benassai2026}.
    \end{itemize}
    \item \textbf{Misure di contorno o a opportunità cumulative (\textit{Contour/Cumulative Opportunities})}: 
        fissano una soglia massima di tempo o distanza e contano il numero complessivo di servizi raggiungibili entro quel determinato confine \cite{geurs2004,mazzulla2024}. Questa metrica è alla base dei modelli urbanistici di prossimità come il paradigma della "Città dei 15 minuti", \textbf{AGGIUNGERE DESCRIZIONE DI COS'È} \cite{wu2025}.
       \begin{itemize}
            \item \textit{Vantaggi}: considerano la quantità di offerta mantenendo un'elevata trasparenza interpretativa \cite{benassai2026}.
            \item \textit{Svantaggio}: usano una soglia arbitraria per poter calcolare la distanza o il tempo: piccole variazioni possono alterare i risultati e le rispettive classifiche di accessibilità \cite{battiston2024}.\\
                Un secondo svantaggio è il cosiddetto ''effetto bordo'' (\textit{edge effect}), cioè l'assenza di un decadimento graduale con la distanza. Ne consegue che due servizi, rispettivamente a 1 minuto e 14 minuti di distanza, assumono lo stesso valore, mentre un terzo servizio con una distanza maggiore di 15 minuti viene escluso dal calcolo \cite{geurs2004,mazzulla2024}.
       \end{itemize}
    
          \item \textbf{Misure potenziali o gravitazionali (\textit{Gravity-based})}: prendono ispirazione dalla legge di gravitazione di Newton \cite{hansen1959}. Superano il limite delle misure cumulative ponderando ciascuna opportunità con una funzione di decadimento continua. In questo modo, più un servizio è distante, minore sarà il suo costributo all'accessibilità della zona di origine.
        \begin{itemize}
            \item \textit{Vantaggi}: sono teoricamente solide poiché riflettono il comportamento spaziale degli utenti ed eliminano l'effetto bordo \cite{geurs2004}.
            \item \textit{Svantaggi}: il valore numerico risultante è meno immediato da interpretare e comunicare. Inoltre, richiedono una delicata calibrazione empirica del parametro di sensibilità alla distanza ($\beta$) basata su dati reali di mobilità \cite{geurs2004, mazzulla2024}.
        \end{itemize}

        \item \textbf{Misure che includono gli effetti di competizione}: risultano indispensabili quando si analizzano servizi primari caratterizzati da una capacità ricettiva limitata (come ospedali, scuole o spazi verdi), in cui un'elevata concentrazione di popolazione può generare congestione dell'offerta\cite{geurs2004, battiston2024}. 
         Tra gli approcci più diffusi rientrano gli indicatori di disponibilità pro-capite \cite{battiston2024}, i modelli probabilistici di scelta (come il modello di Huff applicato alle strutture sanitarie \cite{tariverdi2023}) e la famiglia dei metodi \textit{Two-Step Floating Catchment Area} (2SFCA) \cite{luo2003, mazzulla2024}, 
         che mettono in relazione la capacità di ciascun servizio con la domanda potenziale residente nel suo bacino di utenza.
\end{enumerate}

\subsection{Misure basate sulla persona e spazio-temporali (\textit{Person-based})}
Fondate sulla "Geografia del Tempo" (\textit{Time Geography}), introdotta dal geografo svedese Hägerstrand \cite{hagerstrand1970}, sono misure disaggregate che analizzano l'accessibilità a livello del singolo individuo. Questa misurazione viene eseguita attraverso la costruzione di prismi spazio-temporali (\textit{space-time prisms}) \cite{geurs2004, mazzulla2024}.
Esse valutano la fattibilità degli spostamenti tenendo conto degli impegni del singolo, i tempi di spostamento e gli orari di apertura dei servizi. Sebbene offrano un'elevatissima precisione teorica e permettano di cogliere le disuguaglianze individuali, la loro applicazione su larga scala ha un limite sulla difficoltà e dati individuali su larga scala \cite{geurs2004}.

\subsection{Misure basate sull'utilità e sull'entropia (\textit{Utility-based})}
Queste metriche hanno origine nella teoria microeconomica e nella teoria dell'utilità aleatoria, esse interpretano l'accessibilità come il beneficio economico \cite{geurs2004,mazzulla2024}.
In alcuni contesti, però, possono mancare dati socio-economici e sono state quindi proposte formulazioni basate sull'\textbf{entropia spaziale} \cite{benassai2026}.
L'entropia misura il grado di diversificazione e omogeneità delle varie opportunità attorno al cittadino: un'alta entropia indica la presenza di molteplici alternative di servizi accessibili in termini di percorrenza che offrono libertà di scelta al singolo utente. \cite{benassai2026}.